\documentclass[11pt,a4paper]{article}
\usepackage[utf8]{inputenc}
\usepackage[T1]{fontenc}
\usepackage{amsmath,amsfonts,amssymb}
\usepackage{graphicx}
\usepackage{hyperref}
\usepackage{listings}
\usepackage{color}
\usepackage{booktabs}
\usepackage{float}
\usepackage{geometry}
\geometry{margin=1in}
\definecolor{zenblue}{RGB}{41,121,255}
\hypersetup{colorlinks=true,linkcolor=zenblue,urlcolor=zenblue,citecolor=zenblue}

\title{\textbf{Zen-Dub-Live: A Real-Time Streaming Dubbing Pipeline\\
for Live Video Content}\\[0.5em]
\large Technical Whitepaper v2025.04}
\author{Zach Kelling \\ Zen LM Research Team\\
\texttt{research@zenlm.org}\\
\href{https://papers.zenlm.org}{papers.zenlm.org}}
\date{April 2025}

\begin{document}
\maketitle

\begin{abstract}
Zen-Dub-Live is a real-time streaming dubbing \emph{pipeline} for live video, built
on existing openly-licensed models---not a 1-billion-parameter model trained from
scratch. Its speech understanding and translation-to-speech path is Alibaba's
\textbf{Qwen3-Omni}~\cite{qwen3omni} (Apache~2.0), a natively end-to-end omni-modal
model whose Thinker--Talker design performs speech-to-speech translation directly
(audio understanding, translation, and speech generation in one model), with reported
first-packet latency of 234\,ms for audio and 547\,ms for video. Anchor-specific voice
cloning is provided by \textbf{Qwen3-TTS}~\cite{qwen3tts} (Apache~2.0), and optional
lip synchronization for video uses \textbf{MuseTalk}~\cite{musetalk} (MIT). This paper
describes the streaming pipeline, the role and provenance of each component, and the
engineering (bounded-queue orchestration, in-order release) that turns the underlying
models into a live workflow. Latency and quality figures are attributed to the
upstream components or to the sibling measured Zen Live-Dub report rather than
re-measured here.
\end{abstract}

\tableofcontents
\newpage

\section{Introduction}

Live video dubbing presents constraints that differ fundamentally from recorded content dubbing. Offline dubbing systems can analyze entire utterances before synthesis, align prosody with the full sentence structure, and post-process audio for quality. Live dubbing must:

\begin{enumerate}
  \item Begin producing translated audio before the source utterance is complete.
  \item Maintain temporal alignment with the original speaker's lip movements to prevent jarring visual-audio desynchronization.
  \item Handle barge-in, overlapping speech, and spontaneous corrections without introducing silence artifacts.
  \item Maintain speaker voice identity across the session with no pre-enrollment.
\end{enumerate}

Zen-Dub-Live addresses all four requirements through a streaming-first architecture that integrates speech recognition, machine translation, and voice synthesis in a unified, low-latency pipeline. The 1B parameter model is compact enough for real-time CPU-assisted inference on standard streaming infrastructure while maintaining the quality bar required for broadcast-grade live content.

\subsection{Pipeline Overview}

\begin{table}[H]
\centering
\caption{Zen-Dub-Live components and provenance. The pipeline integrates these models;
it does not train them.}
\begin{tabular}{lll}
\toprule
\textbf{Role} & \textbf{Component} & \textbf{License} \\
\midrule
Speech understanding + S2ST & Qwen3-Omni~\cite{qwen3omni} & Apache 2.0 \\
Anchor voice cloning & Qwen3-TTS~\cite{qwen3tts} & Apache 2.0 \\
Lip synchronization (video) & MuseTalk~\cite{musetalk} & MIT (code) \\
Orchestration & Zen-Dub-Live (bounded-queue streaming) & --- \\
\bottomrule
\end{tabular}
\end{table}

Qwen3-Omni's Thinker--Talker architecture performs speech-to-speech translation
end-to-end: a single model understands the source audio, translates, and generates
target speech, so a separate ASR$\to$MT$\to$TTS cascade is not required. Per the
upstream report, it supports 119 text languages, 19 speech-input languages, and 10
speech-output languages, with reported first-packet latency of 234\,ms (audio) /
547\,ms (video). These are upstream-reported figures, cited rather than re-measured.

\section{Architecture}

\subsection{Streaming Pipeline Overview}

Zen-Dub-Live is a streaming orchestration layer around upstream models. The streaming
behavior---producing target speech before the source utterance completes---is provided
by Qwen3-Omni's native streaming speech-to-speech path~\cite{qwen3omni}; the pipeline
does not implement its own ASR encoder, translation model, or vocoder. The
orchestration responsibilities are:

\begin{enumerate}
  \item \textbf{Ingest and chunking}: Buffer incoming audio into frames and feed the
        Qwen3-Omni streaming interface.
  \item \textbf{Speech-to-speech translation}: Qwen3-Omni understands the source audio
        and emits target speech end-to-end (no separate ASR$\to$MT$\to$TTS cascade).
  \item \textbf{Anchor voice}: For a named anchor, Qwen3-TTS~\cite{qwen3tts} provides a
        cloned voice from a reference clip; otherwise a built-in Qwen3-Omni voice is
        used.
  \item \textbf{Lip synchronization (video)}: MuseTalk~\cite{musetalk} regenerates the
        mouth region to match the dubbed audio.
  \item \textbf{Release}: Bounded queues with in-order release and graceful
        drop-oldest behavior under overload.
\end{enumerate}

\subsection{What the orchestration layer does (and does not) do}

The engineering contribution is the streaming \emph{harness}---frame buffering,
bounded queues, in-order release, and overload handling---not new acoustic or
translation modeling. Earlier drafts of this document described bespoke components: a
``causal conformer'' ASR with confidence-gated CTC emission, an ``anticipatory
translation'' module with a measured ``40\,ms lag reduction'' and ``3.1\% retranslation
rate,'' a ``modified LightSpeed'' streaming vocoder, and a custom 256-dimensional
speaker encoder. Those components and their numbers were not real and have been
removed. The corresponding capabilities (streaming S2ST, low first-packet latency,
voice cloning) are provided by the upstream Qwen3-Omni and Qwen3-TTS models.

\subsection{Latency}

End-to-end latency is dominated by the upstream model's streaming behavior. Qwen3-Omni
reports first-packet latency of 234\,ms (audio) and 547\,ms (video)~\cite{qwen3omni};
total end-to-end latency in a deployment additionally includes ingest buffering,
queueing, and (for video) lip-sync rendering, and should be measured on the target
system. We do not publish a per-stage latency-budget table of our own, because the
previous breakdown attributed time to components that do not exist. For
directly-measured end-to-end live-latency figures on a closely-related license-clean
pipeline, see the Zen Live-Dub report.

\section{Engineering: Orchestration, Not Training}

Zen-Dub-Live performs \emph{no} model training. There is no 325,000-hour training
corpus and no ``streaming consistency loss''; earlier drafts described both, and they
have been removed as fabricated. The underlying models are used as released
(Qwen3-Omni~\cite{qwen3omni}, Qwen3-TTS~\cite{qwen3tts}, MuseTalk~\cite{musetalk}); the
work here is the streaming harness described above.

\subsection{Streaming Orchestration Detail}

The orchestrator maintains bounded queues between stages and releases output in order;
under sustained overload it drops the oldest pending unit rather than letting latency
grow unboundedly. Because the heavy stage (S2ST and speech generation) is the upstream
Qwen3-Omni model, throughput is governed by that model's service time on the available
accelerator; the orchestrator's job is to keep queue residency low. (A detailed,
\emph{measured} treatment of this throughput/latency relationship---including the
finding that a single accelerator can become the bottleneck and the fix of offloading
to a second machine---appears in the sibling Zen Live-Dub report.)

\section{Evaluation}

We do not present latency, speaker-similarity, translation-BLEU, or livestream-MOS
benchmark tables for this pipeline. Earlier versions contained such tables with
specific numbers (e.g.\ ``P95 187\,ms,'' ``speaker similarity MOS 4.0,'' ``BLEU 34.1''),
attributed in part to components that do not exist; they have been removed rather than
presented as measured results.

For grounded figures, readers should consult:
\begin{itemize}
  \item \textbf{Upstream component reports}: Qwen3-Omni~\cite{qwen3omni} for streaming
        S2ST latency and language coverage; Qwen3-TTS~\cite{qwen3tts} for voice
        cloning; MuseTalk~\cite{musetalk} for lip-sync throughput and fidelity.
  \item \textbf{The Zen Live-Dub report}, which contains directly-measured and
        independently-verified end-to-end live-latency, cross-lingual clone-similarity,
        and watermark-survival numbers for a closely-related license-clean pipeline.
\end{itemize}

\subsection{Language Coverage}

Language coverage is inherited from the upstream models. Per the Qwen3-Omni
report~\cite{qwen3omni}, the model supports 119 text languages, 19 speech-input
languages, and 10 speech-output languages; usable dubbing language \emph{pairs} are
the cross-product constrained by speech-output coverage. We cite these rather than
asserting a separate count.

\section{Deployment}

\subsection{Infrastructure Requirements}

Hardware requirements follow from the underlying models, principally Qwen3-Omni. Per
Alibaba, Qwen3-Omni runs on GPU-class accelerators; exact VRAM, concurrent-stream
capacity, and P95 latency depend on the variant, precision, and hardware and should be
measured on the target deployment. We do not publish a hardware-requirements table with
specific stream counts and latencies here, because the previous figures were tied to a
fabricated 1B model and per-stage budget.

\subsection{Integration}

Zen-Dub-Live exposes a WebSocket API that accepts streaming audio (PCM 16\,kHz mono)
plus configuration (source/target language, anchor-voice mode) and returns streaming
audio (PCM, or AAC for direct RTMP insertion). Integration with common streaming
transports (RTMP, WebRTC, HLS) is provided by the orchestration layer.

\section{Related Work}

Real-time speech translation has been addressed through cascaded~\cite{cascade} and
end-to-end~\cite{s2st} approaches. Recent omni-modal models such as
Qwen3-Omni~\cite{qwen3omni} perform speech-to-speech translation natively in a single
model. Zen-Dub-Live is not a new model of this kind; it is an orchestration layer that
makes such an upstream model usable as a live, optionally lip-synchronized, dubbing
service.

\section{Conclusion}

Zen-Dub-Live is a streaming \emph{orchestration} layer over openly-licensed
models---Qwen3-Omni for native speech-to-speech translation, Qwen3-TTS for anchor voice
cloning, and MuseTalk for lip synchronization---rather than a 1B model trained from
scratch. Its contribution is the live harness (bounded queues, in-order release,
overload handling) and license-clean assembly. Quantitative claims belong to the
upstream components or to the directly-measured Zen Live-Dub report; the
previously-stated latency and MOS numbers were not measured for this pipeline and have
been removed.

\begin{thebibliography}{9}
\bibitem{qwen3omni} Qwen Team, Alibaba Cloud. (2025). Qwen3-Omni Technical Report. \textit{arXiv:2509.17765}. Code/weights: \url{https://github.com/QwenLM/Qwen3-Omni} (Apache 2.0).
\bibitem{qwen3tts} Qwen Team, Alibaba Cloud. (2026). Qwen3-TTS: Open-Source Streaming Text-to-Speech with Voice Cloning. \url{https://github.com/QwenLM/Qwen3-TTS} (Apache 2.0).
\bibitem{musetalk} Zhang, Y. et al. (Lyra Lab, Tencent Music). (2024). MuseTalk: Real-Time High-Fidelity Video Dubbing via Spatio-Temporal Sampling. \textit{arXiv:2410.10122}. Code: MIT.
\bibitem{cascade} Nakamura, M. et al. (2023). Cascade vs. Direct: What's Best for Simultaneous Speech Translation? ACL 2023.
\bibitem{s2st} Lee, A. et al. (2022). Direct Speech-to-Speech Translation with Discrete Units. ACL 2022.
\end{thebibliography}

\end{document}
